% ====================================================================
% VT2 Results -- all tables, generated by
% evaluation/report/build_report.py. Do not edit by hand;
% re-run the script instead.
% ====================================================================

% ---- t1_dataset --------------------------------------------------
\begin{table}[h]
\centering
\caption{Trials in the D23 evaluation campaign, after the campaign filter. \enquote{Compl.} counts trials that ran all $H=400$ steps; the remainder aborted (protective stop, RTDE error, mocap loss or hang watchdog). Sim rollouts mirror only completed real trials.}
\label{tab:dataset}
\begin{tabular}{lrrrrr}
\hline
\textbf{Configuration} & \textbf{Real 3\,cm} & \textbf{Compl.} & \textbf{Real 4\,cm} & \textbf{Compl.} & \textbf{Sim} \\
\hline
L0 (none) & 21 & 14 & 6 & 6 & 20 \\
L1 (pos) & 21 & 17 & 6 & 0 & 17 \\
L2 (pos+cube) & 21 & 21 & 6 & 2 & 23 \\
L3 (pos+cube+robot) & 21 & 13 & 6 & 0 & 13 \\
L4 (full) & 21 & 15 & 6 & 3 & 18 \\
\textbf{Total} & \textbf{105} & \textbf{80} & \textbf{30} & \textbf{11} & \textbf{91} \\
\hline
\end{tabular}
\end{table}

% ---- t2_retention ------------------------------------------------
\begin{table}[h]
\centering
\caption{Return retention on the 3~cm (in-distribution) cube. $R$ is the mean episode return, summed over the exact-parity reward subset only (\texttt{gripper\_box}, \texttt{approach\_open}, \texttt{gripper\_align}, \texttt{grasp}, \texttt{lift}, \texttt{action\_rate}), over paired episodes; retention is $R_{\mathrm{real}}/R_{\mathrm{sim}}$ on raw returns with a paired bootstrap resampling episodes. \enquote{Above floor} is the D9 reportability gate: yes only if the 95\,\% interval of the gap lies entirely outside the $\pm$noise-floor band.}
\label{tab:retention}
\begin{tabular}{lrrrrcrrc}
\hline
\textbf{Configuration} & \textbf{$n_{\mathrm{ep}}$} & \textbf{$R_{\mathrm{sim}}$} & \textbf{$R_{\mathrm{real}}$} & \textbf{Retention} & \textbf{95\,\% CI} & \textbf{Gap} & \textbf{Floor} & \textbf{Above floor} \\
\hline
L0 (none) & 5 & 370 & 616 & 1.66 & [1.00, 2.90] & -246 & 172 & no \\
L1 (pos) & 6 & 3,586 & 1,795 & 0.50 & [0.41, 0.65] & 1,791 & 762 & yes \\
L2 (pos+cube) & 7 & 3,163 & 2,276 & 0.72 & [0.55, 1.06] & 887 & 959 & no \\
L3 (pos+cube+robot) & 5 & 3,524 & 1,777 & 0.50 & [0.37, 0.60] & 1,747 & 576 & yes \\
L4 (full) & 6 & 2,839 & 1,705 & 0.60 & [0.45, 0.76] & 1,134 & 403 & yes \\
\hline
\end{tabular}
\end{table}

% ---- t3_stage_retention ------------------------------------------
\begin{table}[h]
\centering
\caption{Per-stage return retention $R_{\mathrm{real}}/R_{\mathrm{sim}}$ on the 3~cm cube, reward terms grouped by task stage. A dash marks a stage whose simulated return is identically zero, so the ratio is undefined: $L0$ never grasps in either domain.}
\label{tab:stage_retention}
\begin{tabular}{lrrrrr}
\hline
\textbf{Stage} & \textbf{L0 (none)} & \textbf{L1 (pos)} & \textbf{L2 (pos+cube)} & \textbf{L3 (pos+cube+robot)} & \textbf{L4 (full)} \\
\hline
Approach & 1.26 & 0.61 & 0.81 & 0.60 & 0.81 \\
Grasp & -- & 0.08 & 0.21 & 0.12 & 0.06 \\
Lift & -- & 0.20 & 0.72 & 0.15 & 0.03 \\
Transport & -- & 0.14 & 0.34 & 0.10 & 0.01 \\
Regularizer & 1.42 & 1.28 & 1.54 & 1.38 & 1.38 \\
\hline
\end{tabular}

\vspace{2pt}
\begin{minipage}{0.95\linewidth}
\footnotesize Terms per stage: approach = \texttt{gripper\_box}, \texttt{approach\_open}, \texttt{gripper\_align}; grasp = \texttt{grasp}; lift = \texttt{lift}; transport = \texttt{box\_target}, \texttt{hold\_target}; regularizer = \texttt{no\_floor\_collision}, \texttt{robot\_target\_qpos}, \texttt{action\_rate}.
\end{minipage}
\end{table}

% ---- t4_abort ----------------------------------------------------
\begin{table}[h]
\centering
\caption{Aborted real trials per configuration and cube, as aborted/attempted. An abort is a protective stop, RTDE error, mocap loss or hang-watchdog trip; simulation never aborts. $L0$ never closes the gripper, so it cannot fail a grasp -- its 4\,cm column is not evidence of robustness.}
\label{tab:abort}
\begin{tabular}{lcc}
\hline
\textbf{Configuration} & \textbf{3\,cm cube} & \textbf{4\,cm cube} \\
\hline
L0 (none) & 7/21 & 0/6 \\
L1 (pos) & 4/21 & 6/6 \\
L2 (pos+cube) & 0/21 & 4/6 \\
L3 (pos+cube+robot) & 8/21 & 6/6 \\
L4 (full) & 6/21 & 3/6 \\
\hline
\end{tabular}

\vspace{2pt}
\begin{minipage}{0.95\linewidth}
\footnotesize Jaw opening is 49.9~mm (\texttt{mujoco\_playground/\_src/manipulation/my\_ur3/universal\_robots\_ur3e/ur3e\_position.xml}). The 3\,cm cube leaves 10.0~mm per side and its diagonal (42.4~mm) still fits, so every yaw is graspable. The 4\,cm cube leaves only 4.9~mm per side and its diagonal (56.6~mm) exceeds the opening, so the jaws can only close over $\pm17^\circ$ of yaw about each $90^\circ$ symmetry, i.e.\ 38\,\% of orientations.
\end{minipage}
\end{table}

% ---- t5_replay ---------------------------------------------------
\begin{table}[h]
\centering
\caption{Dynamics-replay fidelity over completed real trials. One-step error re-synchronises the model to the measured state every control tick and integrates a single tick; open-loop sets the measured initial state once and then replays the whole logged command sequence without correction. Time-to-divergence is the first tick at which open-loop tool-centre-point error exceeds 1\,cm.}
\label{tab:replay}
\begin{tabular}{lrrrrrr}
\hline
\textbf{Configuration} & \textbf{$n$} & \textbf{One-step TCP (mm)} & \textbf{$E_{\mathrm{replay}}$ (rad$^2$)} & \textbf{TTD (steps)} & \textbf{TTD (s)} & \textbf{Never div.} \\
\hline
L0 (none) & 20 & 1.02 & 7.80e-06 & 15 & 0.30 & 3 \\
L1 (pos) & 17 & 1.14 & 1.01e-05 & 21 & 0.42 & 0 \\
L2 (pos+cube) & 23 & 1.16 & 8.15e-06 & 20 & 0.40 & 0 \\
L3 (pos+cube+robot) & 13 & 0.89 & 7.02e-06 & 76 & 1.50 & 1 \\
L4 (full) & 18 & 0.96 & 7.32e-06 & 15 & 0.30 & 0 \\
\hline
\end{tabular}

\vspace{2pt}
\begin{minipage}{0.95\linewidth}
\footnotesize Replay is driven by the logged servoJ targets, so the policy is not in the feedback loop; it measures model-versus-robot error along the states that policy visited, and is therefore not a policy-independent quantity. Time-to-divergence is not comparable across configurations, since each policy generates a different trajectory.
\end{minipage}
\end{table}

% ---- t6_place_error ----------------------------------------------
\begin{table}[h]
\centering
\caption{Place error on the 3\,cm cube for completed trials: distance from where the cube landed after the scripted gripper release to the centre of the $15 \times 15$\,cm target square (half-width 7.5\,cm). Recomputed from the logged landing position with the base-frame calibration applied; the value written by the campaign script is in the wrong frame and is not used.}
\label{tab:place_error}
\begin{tabular}{lrrrr}
\hline
\textbf{Configuration} & \textbf{$n$} & \textbf{In square} & \textbf{Median err.\ (cm)} & \textbf{Best (cm)} \\
\hline
L0 (none) & 21 & 0 & 27.4 & 26.2 \\
L1 (pos) & 17 & 0 & 27.9 & 25.2 \\
L2 (pos+cube) & 21 & 3 & 40.3 & 3.9 \\
L3 (pos+cube+robot) & 13 & 0 & 30.3 & 26.5 \\
L4 (full) & 15 & 2 & 27.8 & 5.3 \\
\hline
\end{tabular}

\vspace{2pt}
\begin{minipage}{0.95\linewidth}
\footnotesize Real-only: MJX has no scripted release event, so there is no simulated counterpart and this is never folded into the reward-based metrics.
\end{minipage}
\end{table}

% ---- t7_regret ---------------------------------------------------
\begin{table}[h]
\centering
\caption{Sim-selection regret on the 3~cm cube: the real return given up by choosing a configuration on simulated return alone. \enquote{Picked} is how often the bootstrap selects that configuration as the simulated best.}
\label{tab:regret}
\begin{tabular}{lrrr}
\hline
\textbf{Configuration} & \textbf{$R_{\mathrm{sim}}$} & \textbf{$R_{\mathrm{real}}$} & \textbf{Picked} \\
\hline
L2 (pos+cube) & 2,938 & 1,853 & 64\,\% \\
L3 (pos+cube+robot) & 2,777 & 1,072 & 36\,\% \\
L1 (pos) & 2,563 & 1,509 & 0\,\% \\
L4 (full) & 2,390 & 1,615 & 0\,\% \\
L0 (none) & 385 & 490 & 0\,\% \\
\hline
\end{tabular}

\vspace{2pt}
\begin{minipage}{0.95\linewidth}
\footnotesize Regret 0 (95\,\% CI [0, 1025]) over the 3 episodes common to all five configurations in both domains. The point estimate is zero because simulation and reality agree on the best configuration here, but with only 3 shared episodes the interval still admits a substantial loss.
\end{minipage}
\end{table}

% ---- t8_sim_success ----------------------------------------------
\begin{table}[h]
\centering
\caption{Task outcome in simulation on the 3~cm cube: the environment's own success flag over the matched-initialisation rollouts, and the median final cube-to-target distance. Even in simulation the policies succeed on a minority of episodes, so the retention ratios in Table~\ref{tab:retention} compare two imperfect systems rather than a solved simulation against a failing robot.}
\label{tab:sim_success}
\begin{tabular}{lrrrr}
\hline
\textbf{Configuration} & \textbf{$n$} & \textbf{Successes} & \textbf{Rate} & \textbf{Median final dist.\ (cm)} \\
\hline
L0 (none) & 21 & 0 & 0\,\% & 35.4 \\
L1 (pos) & 18 & 2 & 11\,\% & 2.4 \\
L2 (pos+cube) & 23 & 5 & 22\,\% & 12.7 \\
L3 (pos+cube+robot) & 15 & 2 & 13\,\% & 6.9 \\
L4 (full) & 15 & 0 & 0\,\% & 19.9 \\
\hline
\end{tabular}
\end{table}

% ---- t9_coverage -------------------------------------------------
\begin{table}[h]
\centering
\caption{Board conditions against the pose ranges seen during training. \enquote{in} means the condition falls inside the randomization range that configuration was trained with. $L0$ is out of distribution everywhere by construction: the ladder generator collapses every pose parameter to a single point for it, so it saw one cube position, one yaw, a flat table at nominal height and one arm pose.}
\label{tab:coverage}
\begin{tabular}{llllcc}
\hline
\textbf{Cond.} & \textbf{Ep.} & \textbf{What varies} & \textbf{Trained range} & \textbf{$L1$--$L4$} & \textbf{$L0$} \\
\hline
A1--A3 & 0--2 & Field position (P1/P2/P3) & \texttt{box\_xy\_jitter} (0.17, 0.24) & in & \textbf{out} \\
B1--B2 & 3--4 & Cube yaw $45^\circ$ & \texttt{box\_z\_rot\_range} $6.28$ rad (full) & in & \textbf{out} \\
C1 & 5 & Table tilt 5.5$^\circ$ & \texttt{lifter\_tilt\_max} 2.9$^\circ$ & \textbf{out} & \textbf{out} \\
D1 & 6 & Support height 150~mm & \texttt{lifter\_height\_abs} 0.00--0.22~m & in & \textbf{out} \\
E1--E2 & 7--8 & 4~cm cube & \texttt{cube\_size\_xy} $\leq0.020$~m half-extent & \textbf{out} & \textbf{out} \\
\hline
\end{tabular}

\vspace{2pt}
\begin{minipage}{0.95\linewidth}
\footnotesize Ranges are read from \texttt{ur3\_pick.default\_config()} and \texttt{batch\_runs/sweeps/gen\_dr\_ladder.py}; the realized tilt and height are read from the sim mirror's metadata, which reproduces the physically measured condition. The paired comparison remains valid for out-of-distribution conditions, since simulation is given the same condition as the robot; what they do not support is reading the absolute returns as an achievable ceiling.
\end{minipage}
\end{table}

% ---- t10_condition_split -----------------------------------------
\begin{table}[h]
\centering
\caption{Retention split by condition group on the 3~cm cube: the baseline conditions A1--A3, which vary only the field position, against the perturbed conditions B1, B2, C1 and D1. For the four configurations that learn the task the two agree within the sampling error, so the comparison between configurations is not an artifact of the perturbed conditions. $L0$ has a single paired baseline episode and its split is not interpretable.}
\label{tab:condition_split}
\begin{tabular}{lrrrr}
\hline
\textbf{Configuration} & \textbf{$n$ base} & \textbf{Retention base} & \textbf{$n$ pert.} & \textbf{Retention pert.} \\
\hline
L0 (none) & 1 & 3.89 & 4 & 1.22 \\
L1 (pos) & 2 & 0.49 & 4 & 0.50 \\
L2 (pos+cube) & 3 & 0.84 & 4 & 0.64 \\
L3 (pos+cube+robot) & 1 & 0.58 & 4 & 0.48 \\
L4 (full) & 3 & 0.55 & 3 & 0.68 \\
\hline
\end{tabular}

\vspace{2pt}
\begin{minipage}{0.95\linewidth}
\footnotesize Most configurations contribute only one to three baseline episodes, so the retention headline in Table~\ref{tab:retention} is dominated by the perturbed conditions.
\end{minipage}
\end{table}
